% from ega4-8-fr.tex
\oldpage[IV-3]{8}
et comme le foncteur $\varinjlim$ commute au produit tensoriel,
$m=\varinjlim m_\lambda$ est un homomorphisme
$\mathscr A\otimes\mathscr A\to\mathscr A$ de
$\mathcal O_{S_0}$-Modules~; par passage à la limite sur les diagrammes
commutatifs exprimant l'associativité de $m_\lambda$ et l'existence de la
section unité dans $\mathscr A_\lambda$, on voit que $m$ définit sur
$\mathscr A$ une structure de $\mathcal O_{S_0}$-Algèbre et que
$\varphi_\lambda$ est un homomorphisme de $\mathcal O_{S_0}$-Algèbres pour
tout $\lambda\in L$. En outre $\mathscr A$ est limite inductive du système
$(\mathscr A_\lambda,\varphi_{\mu\lambda})$ \emph{dans la catégorie des
$\mathcal O_{S_0}$-Algèbres}, autrement dit, pour toute
$\mathcal O_{S_0}$-Algèbre $\mathscr B$, l'application canonique
\begin{equation}
\operatorname{Hom}_{\mathcal O_{S_0}\text{-Alg.}}(\mathscr A,\mathscr B)
\longrightarrow
\varprojlim_\lambda
\operatorname{Hom}_{\mathcal O_{S_0}\text{-Alg.}}
(\mathscr A_\lambda,\mathscr B)
\tag{8.2.1.1}\label{IV.8.2.1.1-fr}
\end{equation}
qui à tout homomorphisme $f:\mathscr A\to\mathscr B$ de
$\mathcal O_{S_0}$-Algèbres, fait correspondre la famille
$(f\circ\varphi_\lambda)$, est \emph{bijective}. En effet, on sait déjà
qu'elle est injective et identifie
$\operatorname{Hom}_{\mathcal O_{S_0}\text{-Alg.}}(\mathscr A,\mathscr B)$
à une partie de
$\varprojlim_\lambda
\operatorname{Hom}_{\mathcal O_{S_0}\text{-Mod.}}
(\mathscr A_\lambda,\mathscr B)$~; tout revient à voir que si $(f_\lambda)$
est un système inductif d'homomorphismes de $\mathcal O_{S_0}$-Algèbres,
$f_\lambda:\mathscr A_\lambda\to\mathscr B$, sa limite inductive
$f:\mathscr A\to\mathscr B$, qui est par définition un homomorphisme de
$\mathcal O_{S_0}$-\emph{Modules}, est aussi un homomorphisme de
$\mathcal O_{S_0}$-Algèbres~; mais cela résulte du passage à la limite
inductive dans le diagramme commutatif d'homomorphismes de
$\mathcal O_{S_0}$-Modules
\[
\xymatrix{
  \mathscr A_\lambda\otimes\mathscr A_\lambda
    \ar[r]^{f_\lambda\otimes f_\lambda} \ar[d]_{m_\lambda}
    & \mathscr B\otimes\mathscr B \ar[d] \\
  \mathscr A_\lambda \ar[r]_{f_\lambda} & \mathscr B
}
\]
et du fait que le foncteur $\varinjlim$ commute aux produits tensoriels.

On notera enfin que si les $\mathscr A_\lambda$ sont des
$\mathcal O_{S_0}$-Algèbres commutatives, il en est de même de $\mathscr A$.
\end{env}

\begin{env}[8.2.2]
\label{IV.8.2.2-fr}
Supposons maintenant que $S_0$ soit un \emph{préschéma}, et que les
$\mathscr A_\lambda$ soient des $\mathcal O_{S_0}$-Algèbres (commutatives)
\emph{quasi-cohérentes}~; on sait alors que
$\mathscr A=\varinjlim\mathscr A_\lambda$ est une
$\mathcal O_{S_0}$-Algèbre quasi-cohérente (I, 4.1.1). Désignons par
$S_\lambda$ (resp. $S$) le spectre de la $\mathcal O_{S_0}$-Algèbre
$\mathscr A_\lambda$ (resp. $\mathscr A$) (II, 1.3.1), et soient
$u_{\lambda\mu}:S_\mu\to S_\lambda$ (pour $\lambda\leq\mu$) et
$u_\lambda:S\to S_\lambda$ les $S_0$-morphismes correspondant aux
homomorphismes $\varphi_{\mu\lambda}$ et $\varphi_\lambda$ respectivement
(II, 1.2.7)~; il est clair que $(S_\lambda,u_{\lambda\mu})$ est un
\emph{système projectif} dans la catégorie des $S_0$-préschémas. On notera
que les $u_{\lambda\mu}$ et $u_\lambda$ sont des morphismes \emph{affines}
(II, 1.6.2), donc \emph{quasi-compacts} et \emph{séparés}.
\end{env}

\begin{proposition}[8.2.3]
\label{IV.8.2.3-fr}
Avec les notations de (8.2.2), les morphismes $u_\lambda:S\to S_\lambda$
font de $S$ une limite projective du système projectif
$(S_\lambda,u_{\lambda\mu})$ dans la catégorie des préschémas. En outre, si
$h:S_0\to T$ est un morphisme, faisant de tout $S_0$-préschéma un
$T$-préschéma, $S$ est aussi limite projective du système
$(S_\lambda,u_{\lambda\mu})$ dans la catégorie des $T$-préschémas.
\end{proposition}

Prouvons d'abord la seconde assertion de l'énoncé dans le cas $T=S_0$.
